\documentclass[12pt]{standalone}
\usepackage[inline]{asymptote}
\usepackage[utf8]{inputenc}
\usepackage[T1,T2A]{fontenc}
\usepackage[]{graphicx}
\usepackage[russian]{babel}

\begin{document}
    \centering
\begin{asy}
//[width=.7\linewidth]
import graph;
size(7.5cm,0);


real c=2;
real b=8;


path curv=(-b,c)--(b,c)--(b,-c)--(-b,-c)--cycle;
fill(curv, yellow+opacity(.2));


path curv=(-b,c)--(b,c);
draw(curv,gray);


path curv=(-b,-c)--(b,-c);
draw(curv,gray);


//


// --- СЕКРЕТНАЯ ФОРМУЛА ПЕРЕСЧЕТА ---
// Находим, сколько единиц в 1 bp для вашего масштаба:
real user_unit_per_bp = (2 * b) / (7.5 * cm); 

// Толщина вашей линии равна 2bp. Увеличиваем радиус ровно на 2bp:
real line_thickness = 2 * user_unit_per_bp; 
real d = 4 + 0*line_thickness; 
// ------------------------------------


//real d=4+1/cm;
draw(Arc((0,c),d,0,-180),red+2bp);


path curv=(0,c)--(d*Cos(-45),c+d*Sin(-45));
draw(curv,gray,arrow=Arrow(2mm));
label(Label("$r$",Relative(0.5)),curv,WSW);


label(Label("$\delta$"),(d,c),N);
label(Label("$H$"),(-b,0),E);
label(Label("$H$"),(-b,-2c),E);
dot((0,c),L=Label("$A$",black),N);


//


draw(Arc((0,-3c),d,0,180),blue+2bp);
dot((0,-3c),L=Label("$B$",black),N);


\end{asy}    























\end{document}
